\selectlanguage{hebrew}
\section{פרק 8: ארכיטקטורת סוכן יציב: אחזור, אימות ואופטימיזציה (RAG/Type-safety/Prompt-optimization)}

בניית סוכני AI (Artificial Intelligence) המיועדים לסביבת Production דורשת הקפדה על עקרונות ארכיטקטוניים שיבטיחו יציבות, בקרה, ויכולת תחזוקה. שלושה עמודים מרכזיים מהווים את הבסיס לארכיטקטורה כזו: שכבת האחזור (Retrieval), שכבת האימות (Validation), ושכבת האופטימיזציה (Optimization). שכבות אלו, בשילוב עם טכניקות כמו Retrieval-Augmented Generation (RAG), Type Safety, ו-Prompt Optimization, מאפשרות לנו לבנות מערכות AI חזקות ואמינות.

\subsection{1. שכבת האחזור (Retrieval Layer)}

ליבתה של כל מערכת RAG היא היכולת לאחזר מידע רלוונטי ומדויק ממאגר נתונים חיצוני. שכבת האחזור אחראית על גישה למקורות הידע, בין אם מדובר במסדי נתונים, מסמכים, או API חיצוניים, ועל סינון ודירוג המידע הרלוונטי ביותר להקשר הנוכחי \cite{Reger2026}. מטרתה היא לספק למודל השפה הגדול (LLM) את ההקשר הנדרש ליצירת תגובה מנומקת ומדויקת, ובכך למנוע "הזיות" (hallucinations) ולהגביר את מידת הדיוק.

בחירת שיטת האחזור תלויה באופן משמעותי בסוג הנתונים ובדרישות הביצועים. שיטות נפוצות כוללות אחזור מבוסס דמיון וקטורי (vector similarity search), אחזור מבוסס מילות מפתח (keyword search), או שילוב של שתיהן. הדיוק של שכבת האחזור משפיע ישירות על איכות הפלט הסופי, ולכן השקעה באלגוריתמי דירוג יעילים (ranking algorithms) ובבניית אינדקסים אופטימליים היא קריטית \cite{Reger2026}. היבטים כמו Latency ו-Throughput בשכבת האחזור חייבים להיות מנוטרים באופן שוטף כדי להבטיח תגובה מהירה של הסוכן.

\subsection{2. שכבת האימות (Validation Layer) - Type Safety}

לאחר שהמידע אוחזר, יש צורך לוודא את תקינותו, את רלוונטיותו, ואת התאמתו לפורמט הנדרש. כאן נכנסת לתמונה שכבת האימות, ובפרט עקרונות של Type Safety. Type Safety, במובן הרחב, מבטיח שהנתונים והתקשורת בין רכיבי המערכת עומדים בציפיות מוגדרות מראש, ומונעים שגיאות הנובעות מחוסר התאמה בטיפוסי הנתונים או בערכים \cite{Reger2026}.

בארכיטקטורת סוכני AI, Type Safety יכול להתבטא במספר אופנים:

\begin{itemize}
  \item \textbf{אימות מבנה הנתונים (Data Structure Validation):} וידוא שנתוני הקלט והפלט מה-LLM, וכן הנתונים המאוחזרים, תואמים למבנה מוגדר (למשל, JSON schema).
  \item \textbf{אימות טיפוסי הנתונים (Data Type Validation):} בדיקה שהערכים בתוך הנתונים הם מהסוג הנכון (למשל, מספר, מחרוזת, תאריך).
  \item \textbf{אימות טווח ערכים (Value Range Validation):} קביעת גבולות מותרים לערכים מסוימים (למשל, ציון בין 0 ל-100).
  \item \textbf{אימות לוגי (Logical Validation):} בדיקת קשרים לוגיים בין שדות שונים בנתונים.
\end{itemize}

יישום Type Safety חזק מפחית משמעותית את הסיכוי לשגיאות בלתי צפויות, מקל על Debugging, ומבטיח שהסוכן מתקשר בצורה אמינה עם מערכות אחרות. כלים וספריות מתקדמות יכולים לסייע באכיפת כללים אלו באופן אוטומטי, ובכך להוות קו הגנה קריטי מפני כשלים \cite{Reger2026}.

\subsection{3. שכבת האופטימיזציה (Optimization Layer) - Prompt Optimization}

גם כאשר האחזור והאימות מבוצעים כראוי, איכות התגובה הסופית של ה-LLM תלויה באופן קריטי באיכות ה-Prompt שמועבר אליו. שכבת האופטימיזציה מתמקדת בשיפור מתמיד של ה-Prompts כדי להשיג תוצאות מיטביות. Prompt Optimization הוא תהליך איטרטיבי של ניסוח, בדיקה, ושיפור ההנחיות הניתנות למודל \cite{Reger2026}.

אסטרטגיות נפוצות ב-Prompt Optimization כוללות:

\begin{itemize}
  \item \textbf{הוספת דוגמאות (Few-shot prompting):} מתן דוגמאות קונקרטיות של קלט-פלט רצויים כדי להנחות את המודל.
  \item \textbf{הנחיות ברורות ומפורטות:} ניסוח הנחיות חד-משמעיות, תוך ציון כל הדרישות והאילוצים.
  \item \textbf{שימוש ב-Chain-of-Thought (CoT) prompting:} בקשה מהמודל "לחשוב בקול רם" ולפרט את שלבי החשיבה שלו, מה שלרוב משפר את איכות הפתרון.
  \item \textbf{התאמת טמפרטורה (Temperature) ופרמטרים נוספים:} שליטה על מידת האקראיות והיצירתיות של הפלט.
\end{itemize}

תהליך זה משתלב לרוב עם observability (יכולת תצפית) על ביצועי המודל, המאפשרת לזהות Prompts שמובילים לתוצאות פחות טובות ולבצע בהם שיפורים \cite{Reger2026}. אופטימיזציה מתמשכת זו חיונית כדי לשמור על רלוונטיות ואיכות התגובות של הסוכן לאורך זמן, במיוחד כאשר דרישות המשתמש או מאגרי הידע משתנים.

שילוב הדוק של שכבות האחזור, האימות (Type Safety), והאופטימיזציה (Prompt Optimization) יוצר ארכיטקטורה יציבה, אמינה, וניתנת לתחזוקה, המאפשרת בניית סוכני AI אפקטיביים לסביבת Production.
\selectlanguage{english}
